\section{Reusing the R forest package: nuisance provenance and population choice}
\label{exec21:grfr-scope}
A bank considering a retention offer needs an effect for the customers it intends to contact. The R package grf is a useful reference candidate because its wrapper connects nuisance fitting, forest prediction and average-effect estimation. That wrapper is part of the statistical method. A proposed CLI or MCP interface must preserve its population and score construction when exposing the underlying C++ engine. The R package and native core share one implementation lineage; their different deployment routes do not supply independent corroboration. We select the original R package as a bounded reference candidate, with an isolated worker considered for a future prototype. The earlier native-boundary analysis explains the ownership and persistence obligations; here the additional concern is the meaning of supplied nuisance predictions and target options.\citep{exec21_grf_reuse,exec13_grf_memory}

The selected source declares version 2.6.1 at the 30 April 2026 commit. Its requirements include R 3.5.0 or later, C++17, GNU make, Rcpp and RcppEigen, with imports including DiceKriging, lmtest, Matrix and sandwich. A named maintainer and issue tracker establish a maintenance contact route, not a service commitment. The DESCRIPTION declares GPL-3; file-level terms and dependencies need to accompany any distribution decision. The inspected license distinguishes running a program from conveying it. An isolated process is an engineering boundary and does not by itself settle the licensing of a distributed system. These observations justify retaining the original as a reference without claiming a qualified installed environment.\citep{exec21_grf_reuse}

\subsection{What the nuisance inputs must mean}
\label{exec21:grfr-nuisance}
Let $A\in\{0,1\}$ denote the offer, $X$ information available before assignment, and $Y$ the outcome over a fixed horizon. Define $e(x)=P(A=1\mid X=x)$ and $\mu_a(x)=E[Y\mid A=a,X=x]$. Consistency, no interference, conditional exchangeability and positivity give $\tau(x)=\mu_1(x)-\mu_0(x)$ the causal interpretation. These are substantive assumptions about assignment and observation; fitting the forest does not verify them.

The wrapper's outcome nuisance is the marginal regression $m(x)=E[Y\mid X=x]=(1-e(x))\mu_0(x)+e(x)\mu_1(x)$. Missing nuisance vectors are estimated by auxiliary regression forests. Supplied scalar values are replicated, and supplied vectors are checked for length. Such checks cannot establish that a prediction uses only permitted information or that it was estimated without access to the evaluation outcome. The interface therefore needs the feature clock, row identity and fitting provenance as well as the numbers. Initially, the proposed worker retains the package's ordinary nuisance path or an explicitly known randomized propensity; externally fitted nuisance vectors require a separate qualification profile.\citep{exec21_grf_reuse}

Write $m^\star,e^\star,t^\star$ for fixed proposed nuisance functions, with $0<e^\star<1$. The binary score reconstructed by the source is
\begin{equation}
 \Gamma=t^\star+\frac{A-e^\star}{e^\star(1-e^\star)}
       \{Y-m^\star-t^\star(A-e^\star)\}.
 \label{exec21:grfr-score}
\end{equation}
Suppressing $x$, define $u_0^\star=m^\star-e^\star t^\star$ and
$u_1^\star=m^\star+(1-e^\star)t^\star$. Substituting these definitions yields
$\Gamma=u_1^\star-u_0^\star+A(Y-u_1^\star)/e^\star
-(1-A)(Y-u_0^\star)/(1-e^\star)$. For a fresh observation, conditional on these fixed functions, direct expectation gives the project identity
\begin{equation}
 E[\Gamma\mid X]-\tau
 =(e^\star-e)\left\{\frac{u_1^\star-\mu_1}{e^\star}
                   +\frac{u_0^\star-\mu_0}{1-e^\star}\right\}.
 \label{exec21:grfr-bias}
\end{equation}
Indeed, the two residual terms have conditional expectations
$e(\mu_1-u_1^\star)/e^\star$ and
$(1-e)(\mu_0-u_0^\star)/(1-e^\star)$, and collecting errors gives the identity. The conditional bias vanishes when the propensity is correct or when both reconstructed treatment-specific regressions are correct. Correctness of $m^\star$ alone is a different condition. Estimation rates, sample reuse and dependence still govern whether a fitted procedure has valid uncertainty; this fixed-function calculation supplies no finite-sample guarantee.\citep{exec21_grf_reuse,exec21_grf_reference}

Consider a deliberately simple law with no varying $X$: $P(A=1)=1/4$, $Y(0)=0$ and $Y(1)=2$. The true marginal outcome mean is $1/2$ and the effect is two. Supply the correct values $m^\star=1/2$ and $t^\star=2$, but use $e^\star=1/2$. The reconstructed regressions become $-1/2$ and $3/2$, although the true regressions are zero and two. Equation~\eqref{exec21:grfr-score} gives scores one for controls and three for treated observations, so their population mean is $3/2$. Even a correct marginal outcome regression and correct conditional-effect input do not restore the intended mean in this example.

For comparison, the correct propensity $e^\star=1/4$ with deliberately wrong $m^\star=3$ and $t^\star=0$ produces scores four and minus four, whose mean is two. Alternatively, keeping the wrong propensity $1/2$ but setting $m^\star=1$ and $t^\star=2$ reconstructs the correct outcome surfaces and makes both scores two. Exact rational enumeration checks these cases and the bias identity on 108 combinations. These are supplied-input witnesses, not claims that the fitted grf algorithm generates these nuisance values. Their practical purpose is to prevent an adapter from equating a correct outcome model with an unlabeled marginal prediction vector.\citep{exec21_grf_reference}

\subsection{An overlap recommendation changes the target}
\label{exec21:grfr-target}
The function \texttt{average\_treatment\_effect} exposes all, treated, control and overlap populations. Under the preceding causal assumptions, their common representation is
\begin{equation}
 \theta_q=\frac{E[q(X)\tau(X)]}{E[q(X)]},\qquad
 q(X)\in\{1,e(X),1-e(X),e(X)(1-e(X))\}.
 \label{exec21:grfr-populations}
\end{equation}
All denominators must be positive. The source warns when estimated propensities become extreme and suggests alternative targets or filtering. A warning does not prove a failure of identification, establish a universal threshold or authorize changing the requested population. A business owner deciding about all eligible customers must explicitly reconsider the question before receiving an effect for customers with greater assignment overlap.\citep{exec21_grf_reuse}

In a two-group hypothetical bank, let each group contain half the eligible population. Treatment propensities are $1/2$ and $1/10$, and group effects are USD $-4$ and USD 12. The all-customer effect is USD 4. The treated and control population effects are respectively USD $-4/3$ and USD $44/7$. Overlap weights are $1/4$ and $9/100$, giving
\begin{equation}
 \theta_{\rm overlap}
 =\frac{(1/4)(-4)+(9/100)12}{1/4+9/100}
 =\frac{4}{17}.
 \label{exec21:grfr-overlap}
\end{equation}
With an assumed offer cost of USD 2 per recipient, the all-customer margin is USD 2, while the overlap-population margin is USD $-30/17$, approximately USD $-1.76$. Both calculations are correct for their stated populations. Applying the second margin to the original customer population would be a change of decision target, not merely a more stable computation of the same quantity. These dollar figures are reproducible assumptions and do not measure a bank's return.\citep{exec21_grf_reference}

The implementation also makes the method label consequential. For the all-population AIPW path, it aggregates scores and computes centered cluster sums for the standard error. The overlap path instead regresses the outcome residual on the treatment residual, with an intercept and observation weights, and uses a sandwich variance. Inside that branch the AIPW/TMLE method choice is ignored, although the earlier TMLE guard rejects clustered or sample-weighted inputs. A wrapper should report the actual residual-regression method and preserve the combination restrictions. It should not label an overlap result as an all-population augmented estimate merely because the request said AIPW.\citep{exec21_grf_reuse}

The population logic of that residual regression can be checked without executing the package. With correct $m$ and $e$, the treatment residual has conditional mean zero, and
$E[(A-e)(Y-m)\mid X]=e(1-e)\tau(X)$ while
$E[(A-e)^2\mid X]=e(1-e)$. The population regression coefficient therefore equals the overlap effect in Equation~\eqref{exec21:grfr-populations}, provided its denominator is nonzero. This argument explains the target. It does not prove that estimated residuals, a particular sandwich correction or a small sample deliver calibrated intervals.

\subsection{Qualification follows the complete statistical request}
\label{exec21:grfr-contract}
The initial proposal profile is one binary action and complete independent customer rows, with a fixed outcome horizon and all-population AIPW target. Clusters, sample weights, continuous actions, instrumental variables, survival, multiarm comparisons and repeated policies need separate profiles. The source exposes several of these methods, but their presence in a dispatcher does not transfer their assumptions or uncertainty guarantees to the initial bank use. A cohort chosen after observing treatment or outcome also changes the evaluation problem; the documented subset warning recommends selection using baseline features. Adaptive selection using fitted effects needs its own honest evaluation design.\citep{exec21_grf_reuse}

Prediction outputs need equally careful typing. With no \texttt{newdata}, the causal-forest wrapper uses stored or recomputed out-of-bag predictions; explicitly supplying the training matrix as new data follows a different prediction path. Those calls are not interchangeable evidence of sample separation. The documented CATE variance, R-loss diagnostic and Monte Carlo excess error answer different questions. Increasing tree count can address forest simulation instability without creating exchangeability, recovering unsupported customers or resolving dependence. Model serialization must also account for the stored original observations and nuisance vectors.\citep{exec21_grf_reuse}

The future worker should retain target identity, actual estimator, customer keys, column order, units, time horizon, warnings, source/dependency versions and nuisance provenance. It must distinguish unavailable execution, an unsupported request, a refused profile, an audited nonidentification finding and a failed computation. Failed or nonfinite output must not silently publish a partial estimate. Qualification would compare the original package against exact score and aggregation fixtures, test OOB and new-data routing, preserve warnings and refusal behavior, and check persistence and process failures. Public generated laws can then assess bias, interval coverage, widths and failed requests before an expensive bank trial. The small examples above establish the logic of that experiment, not its outcome.

Porting the wrapper or embedding its native core remains an unattractive initial route because it transfers nuisance, targeting, uncertainty and lifetime responsibilities without an observed benefit. Bayesian bcf, XBCF and flexBCF remain substantive comparators with different priors, uncertainty and restoration contracts, as developed in the preceding comparison. An original R worker is therefore the preferred prototype candidate, conditional on its environment and rights review. The measured comparison must include validation, auxiliary fits, training, scoring, serialization, human review and support costs for accepted outputs. No present evidence establishes a speed advantage, lower total cost or bank value for any route.\citep{exec21_grf_reuse,exec21_cxx_bcf,exec21_cxx_xbcf,exec21_cxx_flexbcf}

